\chapter{Research Design and Data Provenance}\label{ch:researchdesign}

\begin{objectives}
Turn an economic idea into a falsifiable, reproducible quantitative study; choose
authoritative data; and preserve point-in-time lineage from source to conclusion.
\end{objectives}

\section{The research question}

A useful question names a population, horizon, information set, target, and loss.
For example: \emph{Does an expanding-window term-structure model improve the
one-month log score of recession probabilities relative to a constant-probability
benchmark, using only releases available at each month end?} This is testable.
\emph{Can AI predict markets?} is not yet a research design.

Before coding, write:

\begin{itemize}
\item economic mechanism and competing explanation;
\item estimand or forecast target;
\item observation and decision unit;
\item sample, universe, and inclusion rules;
\item information-availability and execution timeline;
\item benchmark and loss;
\item planned robustness and failure criteria.
\end{itemize}

\section{A hierarchy of evidence}

Descriptive evidence summarizes a defined sample. Predictive evidence demonstrates
performance on genuinely new observations. Causal evidence requires an intervention
or identification strategy. Structural evidence estimates an economic model under
maintained restrictions. A backtest is not automatically any one of these; its claim
must be named.

Theory disciplines which variables and signs matter. Exploratory visualization finds
errors and hypotheses. A validation set selects models. A locked test estimates
generalization once. Reusing the test turns it into training data.

\section{Primary and secondary data}

Prefer a primary producer for definitions and release metadata:

\begin{longtable}{p{0.19\textwidth}p{0.30\textwidth}p{0.41\textwidth}}
\toprule
Domain & Source & Research use and caution\\
\midrule
U.S. macro & FRED and ALFRED, \url{https://fred.stlouisfed.org/docs/api/fred/} &
API access to observations and vintages; inspect underlying agency, units, seasonal
adjustment, revisions, and release time.\\
U.S. yields & U.S. Treasury, \url{https://home.treasury.gov/resource-center/data-chart-center/interest-rates} &
Official daily par curves and methodology; par yields are interpolated curve values,
not transaction yields or zero rates.\\
Monetary policy & Federal Reserve and New York Fed,
\url{https://www.newyorkfed.org/markets} &
Rates, operations, reference rates, and SOMA material; align business dates and
publication timestamps.\\
U.S. filings & SEC EDGAR APIs, \url{https://www.sec.gov/search-filings/edgar-application-programming-interfaces} &
Company facts and filings; comply with fair-access policy, map identifiers, preserve
filing and acceptance times, and handle amendments.\\
Factors & Kenneth French Data Library,
\url{https://mba.tuck.dartmouth.edu/pages/faculty/ken.french/data_library.html} &
Research factor benchmarks with construction notes; methodology can change and
published portfolios are not a substitute for point-in-time reconstruction.\\
Options & Cboe market statistics,
\url{https://www.cboe.com/data/market_statistics/} &
Official exchange statistics and selected downloads; detailed historical quotes and
trades may require licensed products.\\
Europe & ECB Data Portal, \url{https://data.ecb.europa.eu/} &
Monetary, financial, yield, and macro series with structured metadata; inspect
frequency, seasonal adjustment, and area composition.\\
Global & BIS Data Portal, \url{https://data.bis.org/} &
Credit, banking, liquidity, debt, FX, and property series; cross-country definitions
and breaks require care.\\
\bottomrule
\end{longtable}

Convenience services are useful for prototypes but may change schemas, adjustments,
symbols, and terms. The repository's market-price pipeline uses Yahoo Finance through
\texttt{yfinance} as a convenient adjusted-price feed and labels it as such. Claims
requiring production-quality price history should be repeated with licensed,
point-in-time data.

\section{Data due diligence}

For every field, answer:

\begin{enumerate}
\item Who creates the underlying observation?
\item What exactly is measured and in what units?
\item Is it a level, rate, index, yield, return, or annualized transformation?
\item When did the event occur, when was it published, and when was it revised?
\item What do missing and preliminary values mean?
\item Did methodology, coverage, identifier, or calendar change?
\item May the data be stored and redistributed?
\end{enumerate}

Automated checks include uniqueness, monotonic timestamps, valid domains, missingness,
outliers, cross-field identities, and checksums. A large move is quarantined for
review, not automatically deleted.

\section{Revisions and vintages}

Macro observations often evolve from advance to revised estimates. Let
$x_{t|v}$ be the estimate for economic period $t$ as recorded at vintage $v$.
Forecasts made at $v$ may use only $x_{t|v}$ for releases available then. A
revision study can treat
$x_{t|v_{\mathrm{final}}}-x_{t|v_{\mathrm{initial}}}$ as an outcome and examine
whether revisions are news or noise.

Cross-sectional fundamentals also have point-in-time versions. Filing date, vendor
processing, and restatement history determine availability. Lagging every field by
an arbitrary fixed number of months is a crude defense, not full lineage.

\section{Reproducible outputs}

Each project should emit:

\begin{itemize}
\item a machine-readable configuration;
\item raw or cached data with allowed provenance;
\item validation report;
\item tidy analytical data;
\item model estimates and predictions;
\item tables and figures generated from saved results;
\item an interpretation and model-risk memo;
\item a command that rebuilds everything.
\end{itemize}

\section{Problems}

\begin{exercise}[Core: provenance]
Choose one FRED series and write a complete data contract including underlying
producer, units, frequency, seasonal adjustment, vintage, missing codes, and release
availability.
\end{exercise}
\begin{exercise}[Core: test protocol]
Pre-register a market forecast including all choices that could otherwise be tuned
after observing test returns.
\end{exercise}
\begin{exercise}[Advanced: revision triangle]
Design a database table for real-time macro vintages and a query that reconstructs
the information set at any historical date.
\end{exercise}


\chapter{Empirical Laboratories}

\begin{objectives}
Integrate the theory in end-to-end projects spanning regime risk, the yield curve,
factors, options, and causal event studies.
\end{objectives}

\section{Laboratory I: cross-asset regime risk}

\subsection*{Question}
How do equity, duration, gold, and volatility exposures behave across latent or
observable volatility regimes, and does volatility targeting improve the full
distribution after costs?

\subsection*{Mathematics}
Let return vector $r_t$ have conditional model
\[
r_t\mid S_t=k\sim t_{\nu_k}(\mu_k,\Sigma_k),\qquad
\Prob(S_t=j\mid S_{t-1}=i)=p_{ij}.
\]
Estimate filtered regime probabilities, not smoothed probabilities, for decisions.
Compare with a transparent observable regime based on lagged market volatility.

Portfolio exposure is
\[
w_t=\min\left(w_{\max},
\frac{\sigma^\star}{\widehat\sigma_{t|t-1}}\right)w_0.
\]
Test whether scaling changes volatility, Expected Shortfall, drawdown, turnover, and
factor exposure. A lower drawdown is not evidence of alpha if average exposure fell.

\subsection*{Protocol}
Use adjusted total-return proxies with explicitly documented limitations. Split by
time. Fit all regime thresholds and covariance models in the expanding training
sample. Charge spread and nonlinear stress costs. Bootstrap blocks of strategy
returns and repeat under alternative start dates.

\subsection*{Failure modes}
Smoothed-state leakage, VIX treated as an investable total-return asset, constant
transaction costs in crises, and synthetic fallback data presented as empirical
evidence are prohibited. The repository labels fallback output in provenance.

\section{Laboratory II: Treasury curve factors and forecasting}

\subsection*{Question}
How much yield-curve variation is described by level, slope, and curvature, and do
Nelson--Siegel dynamics forecast yields or bond excess returns out of sample?

\subsection*{Mathematics}
At each date, fit
\[
y_t(\tau)=\beta_t^\transpose\ell(\tau;\lambda)+\varepsilon_t(\tau)
\]
by weighted least squares. Compare the fitted factors with PCA. Model factors by
VAR(1), random walk, and shrinkage VAR. Map forecast yields into zero-coupon price
changes and duration-scaled returns using an explicit par-to-zero approximation or a
proper bootstrap.

\subsection*{Protocol}
Use official Treasury daily par yields or documented FRED constant-maturity series.
Do not call them zero rates. Freeze $\lambda$ in training or tune it within nested
validation. Forecast recursively. Evaluate yield RMSE by maturity, directional
accuracy, and a bond-pricing loss. Account for the Treasury methodology break in
December 2021.

\subsection*{Extensions}
Estimate a state-space dynamic Nelson--Siegel model; include macro factors; decompose
forecast errors by level, slope, and curvature; test whether accuracy survives
post-2020 and methodology-change subsamples.

\section{Laboratory III: factor replication and robust allocation}

\subsection*{Question}
Can public research factors explain cross-asset portfolio returns, and does
factor-aware covariance shrinkage improve an implementable allocation?

\subsection*{Mathematics}
Estimate
\[
r_t-r_{f,t}\ones=\alpha+Bf_t+\varepsilon_t
\]
with HAC inference. Construct
$\widehat\Sigma=B\widehat\Sigma_fB^\transpose+\widehat D$, shrink it toward a
diagonal or constant-correlation target, and solve constrained minimum variance.

\subsection*{Protocol}
Download factor returns with their current construction notes and archive the file
checksum. Align monthly periods, percentages, and risk-free convention. Use only
past data for factor selection and covariance. Compare equal weight, sample
covariance, factor covariance, and shrinkage after turnover.

\subsection*{Interpretation}
A factor alpha is not skill unless benchmark choice was predeclared and holdings are
implementable. Factor exposure can change across regimes. Report active risk and
drawdown alongside Sharpe.

\section{Laboratory IV: arbitrage-aware option surface}

\subsection*{Question}
Can a cleaned option chain produce a smooth risk-neutral density and stable local
volatility without static arbitrage?

\subsection*{Mathematics}
Infer forwards from put--call parity, filter quotes with executable bounds, and fit
call prices subject to decreasing and convex strike constraints. Recover
\[
f^\Q_T(K)=P(0,T)^{-1}\partial_{KK}C(K,T).
\]
Compare spline, parametric total-variance, and constrained quadratic fits.

\subsection*{Protocol}
Use synchronized bid and ask, not last prices. Record underlying, rates, dividends,
exercise style, multiplier, and timestamps. Fit mid prices while penalizing errors
outside bid--ask. Withhold strikes and maturities. Stress numerical derivatives.

\subsection*{Extensions}
Calibrate Heston; compare price residuals with hedge residuals; study the risk-neutral
tail over time; never interpret it as a physical crash probability without an
explicit risk-premium model.

\section{Laboratory V: causal policy event study}

\subsection*{Question}
What is the response of rates, equities, FX, and credit to an identified monetary
policy surprise rather than to the endogenous policy decision?

\subsection*{Design}
Use a narrow window around scheduled announcements and a surprise measured from an
appropriate high-frequency futures instrument. Separate target-rate and path
components by principal components of the futures curve. Regress intraday asset
changes on surprises, then use local projections for lower-frequency responses.

\subsection*{Threats}
Central-bank information effects, overlapping news, time-zone mismatch, contract
roll, anticipated announcements, heteroskedasticity, and small event count all
matter. Report event-level scatterplots and leave-one-event-out sensitivity.

\section{Problems}

\begin{exercise}[Core: project map]
For each laboratory, list the theorem, estimator, numerical method, validation
design, and decision metric it uses.
\end{exercise}
\begin{exercise}[Advanced: independent replication]
Choose one laboratory and specify an independent data source or market that would
provide the strongest falsification test.
\end{exercise}


\chapter{From Student Project to Research Portfolio}

\begin{objectives}
Communicate advanced quantitative work with clarity, build a coherent portfolio of
evidence, and identify a path from foundations to independent research.
\end{objectives}

\section{What a strong repository demonstrates}

A serious financial-engineering portfolio demonstrates:

\begin{itemize}
\item mathematical fluency through derivations and proofs;
\item statistical discipline through identification and uncertainty;
\item market knowledge through conventions and execution;
\item software quality through tests and reproducibility;
\item judgment through limitations and negative results;
\item communication through readable reports and figures.
\end{itemize}

A large number of shallow notebooks is weaker than a small number of complete
research pipelines. Each flagship project should connect a theory chapter to a data
contract, code, tests, generated results, and a model-risk memo.

\section{Suggested sequence}

\begin{enumerate}
\item \textbf{Foundation.} Implement returns, discounting, Black--Scholes, bond
analytics, optimization, and statistical tests with unit and property tests.
\item \textbf{Replication.} Reproduce a textbook result and one published empirical
result, including differences.
\item \textbf{Extension.} Change an assumption, data source, market, or validation
period for a falsification test.
\item \textbf{Synthesis.} Combine forecasting, portfolio construction, costs, and
risk in an end-to-end decision.
\item \textbf{Original question.} Pre-register a mechanism and test it on
point-in-time data.
\end{enumerate}

\section{Writing a research report}

The report structure is:

\begin{enumerate}
\item abstract with question, method, result, and limitation;
\item economic motivation and falsifiable hypotheses;
\item data provenance and sample construction;
\item mathematical model and identification assumptions;
\item estimator and numerical implementation;
\item predeclared validation and benchmarks;
\item results with uncertainty and economic magnitude;
\item robustness, stress, and negative findings;
\item limitations and next experiment;
\item reproducibility appendix.
\end{enumerate}

Put derivations where a reader can audit them. Put large robustness tables in an
appendix. A figure title states the finding; axes state units; caption states sample
and construction.

\section{Reading research}

For every paper, identify:

\begin{itemize}
\item the claim and counterfactual;
\item the economic mechanism;
\item the key equation and its assumptions;
\item data availability and selection;
\item variation that identifies the parameter;
\item uncertainty and multiple testing;
\item whether implementation resembles a tradable decision;
\item the strongest alternative explanation.
\end{itemize}

Reproduce a central table before extending the paper. Disagreement is useful when it
can be localized to data, timing, estimator, or assumption.

\section{Ethics and epistemic standards}

Do not present synthetic data as observed evidence, revised macro history as a
real-time forecast, midquote returns as executable fills, or a selected backtest as
an untouched test. Do not distribute data beyond its license. Protect credentials
and confidential records.

State uncertainty in language. A model estimate is conditional; a p-value is not
truth; a scenario is not a probability forecast; a risk-neutral density is not a
physical belief. Intellectual honesty is a technical skill.

\section{Mastery checklist}

The reader is ready for independent work when they can:

\begin{itemize}
\item prove core pricing and probability results and state hypotheses;
\item derive an estimator and its limiting uncertainty;
\item implement two independent numerical methods for a benchmark;
\item construct point-in-time datasets and audit timestamps;
\item design a walk-forward or causal validation;
\item decompose portfolio risk and trading costs;
\item explain where a model will fail before seeing the result;
\item write a reproducible report that another researcher can challenge.
\end{itemize}

\section{Capstone}

\begin{exercise}[Capstone research dossier]
Complete one flagship laboratory. Deliver a preregistration, data dictionary,
mathematical appendix, tested code, immutable predictions, cost-aware evaluation,
stress results, model card, and ten-page research paper. Include at least one
deliberate falsification test and one negative finding.
\end{exercise}
